\section{Asymptotics of the cap coefficient}\label{sec:asymptotics}

The physical coefficient in \eqref{eq:physical-coefficient} is a normalized ratio \(W/J\). We compute its logarithmic growth. The product supplies the twist dependence; the leading logarithmic contribution of the determinant will be the same for every twist.

At the homogeneous consecutive roots \(t_j=iy_j\), set
\[
 L_a(y)=\log(\sinh^2y+\sin^2(al)),\qquad g=L_0+L_2-L_1,
 \qquad a_i=N\rho(y_i).
\]
We use \(x,T,E,r,\rho\) from Section~\ref{sec:counting}. Up to additive constants independent of \(\delta\), the product and Jacobian formulas give
\[
 \begin{aligned}
 \log|W|&=\sum_{i<j}g(y_i-y_j)-N\sum_iL_1(y_i),\\
 \log|J|&=\log|\det H|+\sum_i\log a_i.
 \end{aligned}
\]
Here \(H\) is exactly the quantile Jacobian \eqref{eq:quantile-jacobian}. Indeed \(W=\prod V\prod G\), and
\(J=\det\partial_{t_j}\log D_i\) with
\(D_i=\exp(2\pi iT(y_i))\); the remaining factors in this change of coordinates are independent of the twist.

We will use the convolution identity
\begin{equation}\label{eq:force-convolution}
 g'*\rho=L_1'
\end{equation}
with symmetric principal value at the singularity of \(L_0'\).
To verify it, the transforms of the tempered distributions \(L_a''\) are
\[
 \widehat{L_a''}(t)
 =\frac{2\pi t\cosh((4-a)lt)}{\sinh(4lt)},\qquad 0\le a\le4.
\]
For \(0<a<4\), this follows by differentiating in \(a\), using
\(\partial_aL_a=2\pi l p_a\), and integrating from \(a=4\), where
\(L_4''=2\operatorname{sech}^2y\) has the displayed transform by the Beta integral. At \(a=0\) take the tempered distribution limit, which corresponds to symmetric excision for \(L_0'\).
Now
\[
 \cosh(4v)+\cosh(2v)-\cosh(3v)
 =(2\cosh v-1)\cosh(3v).
\]
Multiplication by \(\widehat\rho=(2\cosh v-1)^{-1}\) proves the identity for second derivatives; oddness of the first derivatives fixes the integration constant and yields \eqref{eq:force-convolution}.

\subsection{Bulk limits and the twist derivative}

Write
\[
 \boldsymbol h=H^{-1}\boldsymbol1=\partial_\delta(x_i)_i,
 \qquad \kappa_i=\min(i,N+1-i).
\]
The equality follows by differentiating the counting equations, and
Proposition~\ref{prop:real-counting} bounds \(h_i\) uniformly. We first establish the following meaning of a bulk limit: for every sequence
\(N\to\infty\), \(|\delta|\le1/4\), and \(\kappa_i\to\infty\),
\begin{equation}\label{eq:bulk-counting-limits}
 \int_{\mathbb R}\varphi(z)\bigl(E(y_i+z)+s\bigr)\,dz\longrightarrow0
 \quad(\varphi\in C_c^\infty(\mathbb R)),\qquad
 h_i\longrightarrow\frac43=\frac{ds}{d\delta}.
\end{equation}
The assertion allows \(\delta\) and the row \(i\) to vary with \(N\).

Here is the quantization argument for the first limit. On every fixed rapidity window about a bulk \(y_i\),
\[
 f(y):=T'(y)=N\rho(y)+(r'*E)(y)\asymp a_i\longrightarrow\infty,
 \qquad |f^{(m)}(y)|\le C_m a_i
\]
for fixed \(m\). The correction \(r'*E=r'*(E+s)\) is bounded, and the assertions for \(N\rho\) follow from the density bounds in Section~\ref{sec:counting}. Thus \(T\) is increasing there for large \(N\). The quantile boxes put atoms on both flanks of a slightly smaller window. Let \(y(q)\) be the inverse coordinate for \(q=T(y)-\delta\) on the larger window. The atoms occur at the exact consecutive values \(j-\tfrac12\). On a complete unit cell between successive atoms,
\[
 \#\{j:y_j\le y\}-q=\lfloor q+\tfrac12\rfloor-q
\]
has integral zero in \(q\). For a smooth test function supported in the smaller window, the derivative of its transformed weight is
\[
 \frac{d}{dq}\left(\frac{\varphi(y(q)-y_i)}{f(y(q))}\right)
 =O(a_i^{-2}).
\]
Summing the zero-mean cancellation over \(O(a_i)\) complete cells and bounding the two end cells by \(O(a_i^{-1})\) proves that
\(\#\{j:y_j\le y\}-(T(y)-\delta)\) tends weakly to zero on the window.

Take a weak-star subsequential limit of the bounded translated discrepancies and a limit \(\delta_*\) of the twists. Since
\[
 \#\{j:y_j\le y\}-(T(y)-\delta)=E-r*E+\delta,
\]
the preceding cancellation and the integrable tails of \(r\) show that the limit \(E_*\) on the whole line satisfies
\[
 E_*=r*E_*-\delta_*.
\]
The bounded solution is uniquely the constant \(-s_*=-4\delta_*/3\): subtracting that constant leaves \(U=r*U\), and
\(\|r\|_1<1\) forces \(U=0\). This proves the first assertion of
\eqref{eq:bulk-counting-limits} for every bulk sequence. Kernel decay then upgrades it to
\begin{equation}\label{eq:bulk-smooth-correction}
 (r^{(m)}*E)(y_i+z)\longrightarrow0
 \quad\hbox{locally uniformly in }z,\qquad m\ge1.
\end{equation}
One first restricts the convolution to a compact window and uses weak convergence, then removes the tails using the bound on \(E+s\); the same argument for the next derivative gives local uniformity.

For the limit of \(h_i\), write \(d_j=f_j/a_j\). The linearized equations are
\[
 d_j h_j=1+\sum_k r(y_j-y_k)\frac{h_k}{a_k}.
\]
On a bulk window \(d_j\to1\) uniformly. The sums on the right, viewed as functions of the rapidity in place of \(y_j\), have uniformly bounded derivatives and uniformly negligible tails by
\eqref{eq:counting-kernel-tails}. They therefore have locally uniform subsequential limits. At the atoms, the linearized equation and \(d_j\to1\) make \(h_j\) converge to one plus that continuous limit. The unit quantile cells form a rapidity grid of vanishing mesh in the window, as in the bulk Jacobian estimate. Passing to the limit in the sums gives a bounded continuous function satisfying
\[
 h_*=1+r*h_*.
\]
Its only bounded solution is the constant
\((1-\int r)^{-1}=4/3\), again by \(\|r\|_1<1\). This proves the second assertion of \eqref{eq:bulk-counting-limits}.

Set
\[
 \Phi_N(\delta)=\log|W|-\sum_i\log a_i.
\]
Since \(\partial_\delta y_i=h_i/a_i\), differentiation and
\eqref{eq:force-convolution} give
\begin{equation}\label{eq:product-force}
 \begin{aligned}
 \Phi_N'(\delta)&=\sum_i\frac{h_i}{a_i}B_i,\\
 B_i&=\sum_{j\ne i}g'(y_i-y_j)
       -N\,\operatorname{pv}\!\int_{\mathbb R}g'(y_i-y)\rho(y)\,dy
       -\frac{\rho_i'}{\rho_i}.
 \end{aligned}
\end{equation}
We next prove that \(B_i\) is uniformly bounded and that, along every bulk sequence,
\begin{equation}\label{eq:bulk-force}
 B_i=-4s+o(1).
\end{equation}

Choose an even smooth compactly supported function \(\chi\) equal to one near zero. The singularity of \(g'\) is \(2/y\), so
\[
 g'(y)=b(y)+k(y),\qquad k(y)=\frac{2\chi(y)}y,
\]
where \(b\) is smooth and odd. The explicit formulas for \(L_a\) give
\(g'(y)=\pm2+O(e^{-2|y|})\) and \(g''(y)=O(e^{-2|y|})\) at the two ends. Hence \(b'\) is integrable and \(\int b'=4\). The value \(b(0)=0\) lets us include the self atom in its sum. Integration by parts against the discrepancy gives its entire contribution to the difference in \eqref{eq:product-force}:
\[
 \sum_j b(y_i-y_j)-N\int b(y_i-y)\rho(y)\,dy
 =\int b'(y_i-y)E(y)\,dy.
\]
It is uniformly bounded; by \eqref{eq:bulk-counting-limits} and the integrable tails of \(b'\), it tends to \(-s\int b'=-4s\) in the bulk.

For a fixed edge layer \(\kappa_i\le M\), the singular contribution is also uniformly bounded. On the fixed support window of \(\chi(y_i-\cdot)\), the density \(N\rho\) and its first derivative are bounded by a constant depending on \(M\). The quantile gaps are at least \(0.58\), so the rapidity gaps between atoms in this window have a positive lower bound depending on \(M\); there are only boundedly many such atoms. This bounds the lattice sum. For the principal value integral, subtract the constant \(a_i=N\rho_i\), whose integral against the odd kernel \(k\) is zero, and use
\[
 \left|\operatorname{pv}\!\int k(t)N\rho(y_i-t)\,dt\right|
 \le 2\int\frac{|\chi(t)|}{|t|}
       |N\rho(y_i-t)-N\rho_i|\,dt.
\]
The first-derivative bound makes the right side finite uniformly on the layer. The remaining term \(\rho_i'/\rho_i\) is uniformly bounded by the density estimate.

\paragraph{The missing central value in the singular sum.}
Consider now a bulk sequence. By \eqref{eq:bulk-smooth-correction},
\(f-N\rho=r'*E\to0\) smoothly on the support window. Subtracting its central value inside the odd principal value integral shows
\[
 \operatorname{pv}\!\int k(y_i-y)\bigl(N\rho(y)-f(y)\bigr)\,dy=o(1);
\]
for example, its absolute value is bounded by a constant times the supremum of \(|(N\rho-f)'|\) on that window. We may therefore replace \(N\rho\) by \(f\) in the singular integral at cost \(o(1)\).

On a slightly larger fixed window \(T\) is increasing, and the quantile boxes put atoms on both flanks of the support of \(\chi(y_i-\cdot)\). Let \(y_i(u)\) denote the inverse coordinate
\[
 u=T(y)-T(y_i),\qquad y_i(0)=y_i.
\]
Every atom in the support window has \(u=j-i\). Conversely, every integer in the image of that support lies between the labels of the two flanking atoms. Since the labels are consecutive and \(T\) is increasing, it is realized by exactly one atom in the window. The symmetric principal value in \(y-y_i\) is also the symmetric principal value in \(u\): for fixed \(N\),
\[
 T(y_i\pm\varepsilon)-T(y_i)
 =\pm f_i\varepsilon+\tfrac12f_i'\varepsilon^2+O(\varepsilon^3),
\]
so the ratio of the two absolute excision endpoints tends to one.
Extend
\[
 K_i(u)=\frac{2\chi(y_i-y_i(u))}{y_i-y_i(u)}
\]
by zero outside the inverse window; its support is interior to that window. Choose one even smooth cutoff \(\psi\), equal to one near zero, such that the support of \(\psi(u/f_i)\) is contained in the inverse window for every sufficiently large member of the bulk sequence. This is possible because \(f/f_i\) is bounded above and below on the fixed rapidity window. Subtract the odd principal part by defining, for \(u\ne0\),
\[
 R_i(u)=K_i(u)+\frac{2f_i}{u}\psi(u/f_i).
\]
The sign and the central value follow directly from the inverse Taylor expansion:
\[
 y_i(u)-y_i=\frac{u}{f_i}-\frac{f_i'}{2f_i^3}u^2+O_i(u^3),
 \qquad
 \frac2{y_i-y_i(u)}
 =-\frac{2f_i}{u}-\frac{f_i'}{f_i}+O_i(u).
\]
Because both cutoffs equal one near zero, \(R_i\) extends smoothly there with
\begin{equation}\label{eq:missing-central-value}
 R_i(0)=-\frac{f_i'}{f_i}.
\end{equation}

We spell out the quadrature estimate needed for this central value. Set
\(Y_i(v)=y_i(f_i v)-y_i\). The bounds on \(f/f_i\) and
\(f^{(m)}/f_i\), \(m\le3\), give uniform \(C^4\) bounds for \(Y_i\) on the relevant fixed \(v\)-interval, with \(Y_i'(0)=1\) and \(Y_i(v)/v\) bounded away from zero. The integral Taylor formula for
\((Y_i(v)-v)/v^2\) then shows that \(R_i(f_i v)\) is bounded in
\(C^2\), with support in a fixed compact interval. The midpoint estimate on the cells \([n-\tfrac12,n+\tfrac12]\) gives
\[
 \left|\sum_{n\in\mathbb Z}R_i(n)-\int_{\mathbb R}R_i(u)\,du\right|
 \le\frac1{24}\sum_{n\in\mathbb Z}
       \sup_{|u-n|\le1/2}|R_i''(u)|
 \le\frac C{f_i}.
\]
The last inequality uses \(O(f_i)\) nonzero cells and
\(\|R_i''\|_\infty=O(f_i^{-2})\). The lattice sum and principal value integral of the subtracted odd function
\(-2f_i\psi(u/f_i)/u\) both vanish by symmetry. Exact consecutive quantization now yields
\[
 \begin{split}
 \sum_{j\ne i}k(y_i-y_j)
   -\operatorname{pv}\!\int k(y_i-y)f(y)\,dy
 &=\sum_{n\ne0}R_i(n)-\int R_i(u)\,du\\
 &=-R_i(0)+O(1/f_i)
 =\frac{f_i'}{f_i}+O(1/f_i).
 \end{split}
\]
Finally \eqref{eq:bulk-smooth-correction} implies
\[
 \frac{f_i'}{f_i}
 =\frac{N\rho_i'+(r''*E)(y_i)}
        {N\rho_i+(r'*E)(y_i)}
 =\frac{\rho_i'}{\rho_i}+o(1).
\]
This is precisely the term subtracted in \(B_i\). Combining it with the smooth contribution proves \eqref{eq:bulk-force}. The fixed-layer bound and the bulk limit together also prove
\begin{equation}\label{eq:uniform-force-bound}
 \sup_{\substack{N\ge N_0,\ |\delta|\le1/4\\1\le i\le N}}|B_i|<\infty.
\end{equation}
Indeed, a sequence violating this bound would either have a subsequence in a fixed edge layer, or have \(\kappa_i\to\infty\), and both cases have just been bounded.

\paragraph{Uniform harmonic averaging.}
We give the summation statement that turns these local limits into a uniform twist derivative. Put \(w_i=1/a_i\). The boxes and
\eqref{eq:quantile-density} give \(a_i\asymp\kappa_i\) uniformly. Moreover,
\begin{equation}\label{eq:harmonic-mass}
 \sup_{|\delta|\le1/4}
 \left|\frac1{\log N}\sum_{i=1}^Nw_i-\frac2{\alpha_{\rm r}}\right|
 \longrightarrow0.
\end{equation}
Here is a two-cutoff proof. Fix a fractional cutoff \(0<\eta<1/4\) and a fixed edge cutoff \(M\). The layers \(\kappa_i<M\) have total \(w\)-mass \(O_M(1)\). The middle interval \(\kappa_i>\eta N\) has total mass \(O_\eta(1)\), because \(w_i\le C/(\eta N)\) there. On the left remaining range \(M\le i\le\eta N\), the boxes give \(x_i=i+O(1)\) uniformly in \(\delta\). Given \(\varepsilon>0\), choose \(\eta\) small enough that the left-tail ratio in \eqref{eq:quantile-density} is within \(\varepsilon\) of one whenever \(u/N\le2\eta\), and choose \(M\) large enough that \(x_i/i\) is within \(\varepsilon\) of one for \(i\ge M\). For large \(N\), throughout this range,
\[
 w_i=\frac{1+O(\varepsilon)}{\alpha_{\rm r}i}
\]
uniformly in \(\delta\). The right remaining range has the same estimate with \(N+1-i\). Since
\(\sum_{M\le i\le\eta N}i^{-1}=\log N+O_{M,\eta}(1)\),
first \(N\to\infty\) and then \(\varepsilon\to0\) prove
\eqref{eq:harmonic-mass} uniformly.

We will use the following consequence. Suppose arrays \(z_{N,i}(\delta)\) and a function \(z_*(\delta)\) are uniformly bounded and satisfy
\begin{equation}\label{eq:uniform-bulk-condition}
 \lim_{M\to\infty}\limsup_{N\to\infty}
 \sup_{\substack{|\delta|\le1/4\\ \kappa_i\ge M}}
 |z_{N,i}(\delta)-z_*(\delta)|=0.
\end{equation}
Then
\begin{equation}\label{eq:harmonic-averaging}
 \sup_{|\delta|\le1/4}
 \left|\frac1{\log N}\sum_iw_i z_{N,i}(\delta)
       -\frac2{\alpha_{\rm r}}z_*(\delta)\right|\longrightarrow0.
\end{equation}
To prove it, choose \(M\) so the error in \eqref{eq:uniform-bulk-condition} is small for large \(N\). Its weighted sum outside the fixed edge layers is small times \(\sum_iw_i\); inside those layers it is \(O_M(1)\). Divide by \(\log N\) and use \eqref{eq:harmonic-mass}. This proves the assertion with the supremum in \(\delta\). A scalar bulk limit proved along every sequence with \(N\to\infty\) and \(\kappa_i\to\infty\), allowing arbitrary twists in the compact interval, implies \eqref{eq:uniform-bulk-condition} by contradiction and subsequence extraction.

Apply this rule to \(z_{N,i}=h_iB_i\). It is uniformly bounded by the inverse estimate and \eqref{eq:uniform-force-bound}; by
\eqref{eq:bulk-counting-limits} and \eqref{eq:bulk-force}, its bulk limit is
\(-4s\,ds/d\delta\). Therefore \eqref{eq:product-force} gives
\begin{equation}\label{eq:uniform-product-derivative}
 \frac{\Phi_N'(\delta)}{\log N}
 =-\frac8{\alpha_{\rm r}}s\,\frac{ds}{d\delta}+o(1)
\end{equation}
uniformly for \(|\delta|\le1/4\).

\subsection{The determinant contribution}

It remains to show that the leading logarithmic contribution of \(\det H\) is independent of \(\delta\). Factor
\[
 H=\operatorname{diag}(d_i)(I-A),\qquad
 d_i=\frac{f_i}{a_i},\qquad
 A_{ij}=\frac1{d_i}\frac{r(y_i-y_j)}{a_j}.
\]
The strict margin \eqref{eq:counting-dominance} and the uniform upper bound on \(d_i\) imply
\[
 d_i\ge\epsilon,\qquad
 \sum_j|A_{ij}|\le q_0<1
\]
with absolute \(\epsilon,q_0\), uniformly throughout the rectangles. The diagonal logarithms are negligible at the scale \(\log N\). In fact \(f_i-a_i=(r'*E)(y_i)\) is uniformly bounded, and the logarithm is Lipschitz on the fixed interval containing all \(d_i\). Hence
\[
 |a_i\log d_i|\le C|(r'*E)(y_i)|.
\]
The right side is uniformly bounded and tends to zero in every bulk limit by
\eqref{eq:bulk-smooth-correction}. The averaging rule
\eqref{eq:harmonic-averaging}, applied to \(z_{N,i}=a_i\log d_i\), gives
\begin{equation}\label{eq:diagonal-log-negligible}
 \frac1{\log N}\sum_i\log d_i=o(1)
\end{equation}
uniformly in \(\delta\).

The path \(I-tA\), \(0\le t\le1\), is invertible because
\(\|tA\|_\infty<1\). Its real determinant remains positive, and the convergent trace series is
\[
 \log\det(I-A)=-\sum_{k\ge1}\frac{\operatorname{tr}A^k}{k}.
\]
For all \(N,i,\delta\) and \(k\ge1\), its diagonal terms have the summable bound
\begin{equation}\label{eq:trace-term-bound}
 a_i|(A^k)_{ii}|\le Cq_0^{\,k-1}.
\end{equation}
Indeed \(a_i|A_{ji}|=|r(y_j-y_i)|/d_j\le C\), while
\(\sum_j|(A^{k-1})_{ij}|\le q_0^{k-1}\); multiply these two bounds in the final factor of \(A^k\).

For fixed \(k\), the bulk limit of the same quantity is
\begin{equation}\label{eq:trace-bulk-limit}
 a_i(A^k)_{ii}\longrightarrow (r^{*k})(0).
\end{equation}
For \(k=1\) this is \(r(0)/d_i\to r(0)\). For larger \(k\), expand the product as
\[
 a_i(A^k)_{ii}
 =\sum_{j_1,\ldots,j_{k-1}}
 \frac{r(y_i-y_{j_1})r(y_{j_1}-y_{j_2})\cdots
       r(y_{j_{k-1}}-y_i)}
 {d_i d_{j_1}\cdots d_{j_{k-1}}\,
  a_{j_1}\cdots a_{j_{k-1}}}.
\]
First restrict each rapidity step to a fixed bounded interval. All intermediate atoms then lie in a fixed bulk window, where \(d_j\to1\) uniformly and the weighted quantile grids converge to Lebesgue measure. Iterated quadrature gives the corresponding truncated convolution. The restriction can be removed uniformly: every row of \(|A|\) has negligible mass outside a large rapidity step by
\eqref{eq:counting-kernel-tails}, and the terminal factor
\(a_i|A_{ji}|=|r(y_j-y_i)|/d_j\) is bounded and decays exponentially with its step. For fixed \(k\), these bounds control the union of the finitely many omitted step ranges. This proves
\eqref{eq:trace-bulk-limit} along every bulk sequence.

Apply \eqref{eq:harmonic-averaging} for each fixed \(k\) to
\(z_{N,i}=a_i(A^k)_{ii}\). Equations
\eqref{eq:trace-term-bound}--\eqref{eq:trace-bulk-limit} give
\[
 \frac{\operatorname{tr}A^k}{\log N}
 =\frac2{\alpha_{\rm r}}(r^{*k})(0)+o(1)
\]
uniformly in \(\delta\). Passing this limit through the trace series is justified by the same bound:
\[
 \frac{|\operatorname{tr}A^k|}{\log N}
 \le Cq_0^{k-1}\frac{\sum_iw_i}{\log N}
 \le C' q_0^{k-1}
\]
for all large \(N\), uniformly in \(\delta\). After division by \(k\), the right side is summable. Thus
\[
 \frac{\log\det H}{\log N}
 =-\frac2{\alpha_{\rm r}}\sum_{k\ge1}\frac{(r^{*k})(0)}k+o(1)
\]
uniformly in \(\delta\), where \eqref{eq:diagonal-log-negligible} was used. The series limit is independent of the twist.

Integrate \eqref{eq:uniform-product-derivative} from the vacuum value
\(\delta=1/4\) to \(\delta=\theta/(2\pi)\), and subtract the determinant contributions at those two values. Uniformity of both estimates gives
\[
 \frac{\log\left|(W/J)_\theta/(W/J)_{\pi/2}\right|}{\log N}
 =-\frac4{\alpha_{\rm r}}
       \bigl(s_\theta^2-s_{\pi/2}^2\bigr)+o(1),
 \qquad 0\le\theta\le\frac\pi2.
\]
By \eqref{eq:physical-coefficient} this is the logarithmic asymptotic of
\(c_{\rm phys}(\theta)\). The bulk row has spin twist
\(K^{-1}=e^{-i\theta}\) and winding fugacity
\(K+K^{-1}=2\cos\theta\), while the two vacuum caps remain fixed. Since
\(s_\theta=2\theta/(3\pi)\) and \(\alpha_{\rm r}=8/3\), we obtain
\begin{equation}\label{eq:amplitude-exponent}
 \frac{\log|c_{\rm phys}(\theta)|}{\log N}
 =\frac16-\frac{2\theta^2}{3\pi^2}+o(1),
 \qquad 0\le\theta\le\frac\pi2,
\end{equation}
with the error uniform on this interval. At zero winding fugacity,
\(\theta=\pi/2\), the coefficient is exactly one. At physical fugacity \(2\), \(\theta=0\), it is \(N^{1/6+o(1)}\). The next section uses the geometric interpretation and height estimates to turn this coefficient into a planar nesting law.
